\documentclass[11pt]{article}
\usepackage{fontspec}
\usepackage[french,english]{babel}
\usepackage{amsmath}
\usepackage{amssymb}
\usepackage{amsthm}
\usepackage{mathrsfs}
\usepackage{tikz-cd}
\usepackage{enumitem}

\newcommand{\printedpage}[1]{\marginpar{\footnotesize p.~#1}}
\newcommand{\msp}[1]{\marginpar{\footnotesize [#1]}}
\newcommand{\Hot}{\textsf{Hot}}
\newcommand{\Dia}{\textsf{Dia}}
\newcommand{\Ob}{\textsf{Ob}}
\newcommand{\Hom}{\textsf{Hom}}
\newcommand{\DD}{\mathbb{D}}

\frenchsetup{AutoSpacePunctuation=false}

\begin{document}

\printedpage{59}
% head: A. Grothendieck : Dérivateurs, Ch. I, Édition des Universités de Montpellier II et Paris VII
% foot: Version 8 juillet 2001

\msp{page 96}
Finally, it seems that nothing remains of proposition 2 (p.79) introducing the notion of $\DD$-fibration – not even the equivalence of the two conditions $b$) and $b'$) (universal left $\DD$-fibration, universal right $\DD$-fibration). One thinks of their relation to a condition $a$) modified as follows (for $f : X \to Y$) :

$a$) For every $Y'' \to Y' \to Y$, with $Y'$ and $Y''$ \Hot-\underline{aspheric} (and not only $\DD$-aspheric) $X'' = X \times_Y Y'' \to X' = X \times_Y Y'$ is a $\DD$-equivalence. (When one assumes that $Y''$ and $Y'$ have final elements, this is the condition that $f$ be a universal left $\DD$-fibration; when one assumes that $Y''$ and $Y'$ have initial elements, this is the condition universal right $\DD$-fibration.)\footnote[48]{The reasoning must be made explicit. It suffices to test even on $Y'' \to Y'$, where $Y''$ is \underline{one-point}, and where $Y'$ has a final (resp. initial) object.}

We obviously have $a$) $\Rightarrow$ $b$) and $a$) $\Rightarrow$ $b'$) and we would like these to be equivalences. Here is an interesting axiom\footnote[49]{\textit{cf.} \textbf{Der 3} $a'$) on the following page.} on $\DD$, which seems to hold in all known cases, which ensures the desired equivalence. At the same time I am going to reformulate \textbf{Der 3} in three parts (\textit{cf.} p. 67).

\msp{page 97}
\textbf{Der} 3 $a$) If $Z \in \Dia$ is \Hot-aspheric, then $p : Z \to e$ is a $\DD$-equivalence, \textit{i.e.} for example $p^* : \mathcal{A} = \DD(e) \to \DD(Z)$ is fully faithful.

$b$) Let $f : Z \to Z'$ be an arrow in \Dia. If $f$ is a \Hot-equivalence, it is a $\DD$-equivalence (\textit{i.e.} $f^* : \DD(Z') \to \DD(Z)$ induces an \underline{equivalence} between the subcategories of $Z$, $Z'$ \texttt{[rather of} $\DD(Z)$\texttt{,} $\DD(Z')$\texttt{]} consisting of the constant objects).

$a'$) (Makes $a$ considerably more precise). If $Z \in \Dia$ is \Hot-aspheric, then we have $a$), and the essential image of $\DD(e) = \mathcal{A}$ in $\DD(Z)$ consists of the \underline{locally} \underline{constant} objects of $Z$, \textit{i.e.} of the objects such that the functor
\[
Z^{\circ} \longrightarrow \mathcal{A} \qquad z \longmapsto F_z
\]
is locally constant (\textit{i.e.} sends every arrow to an isomorphism) hence isomorphic to a constant functor (since $Z$ is \Hot-aspheric, hence 1-connected).

It must be clearly noted that in $a'$), it is \underline{not} enough to assume that $Z$ is $\DD$-aspheric, instead of \Hot-aspheric. As soon as $\DD$-equivalence is strictly weaker than \Hot, it seems that $a'$) always becomes false, as a condition on $\DD$, if one assumes there only that $Z$ is $\DD$-aspheric. For example, even if $\DD = \DD_i$ (theory of relative homotopy $i$-types), the category
\msp{page 98}
of locally constant objects on $Z$ is expressed in terms of the homotopy $(i+1)$-type, but not in terms of its homotopy $i$-type alone.

It does seem that it is axiom \textbf{Der} 3 ($a'$) which ensures the existence of a good notion (``bilateral'') of $\DD$-fibration, which is stable under base change, and which reduces either
\printedpage{60}
% head: A. Grothendieck : Dérivateurs, Ch. I, Édition des Universités de Montpellier II et Paris VII
% foot: Version 8 juillet 2001
to ``universal left $\DD$-fibration'' or to ``universal right $\DD$-fibration'' and thus which is in addition a \underline{verifiable} condition in practice. In the case where $f$ is proper or smooth, the criterion is particularly simple, since it is limited to the condition that the $f_* f^* a_Y$, or the $f_! f^* a_Y$ (for $a \in \mathcal{A}$), be locally constant objects on $Y$ (hence \underline{constant} if $Y$ is \underline{\underline{\Hot}}-aspheric, by virtue of \textbf{Der} 3 ($a'$). % [sic] localements -> localement, plural ending on the adverb; % [sic] "(donc ... ($a'$)." -> closing parenthesis after "($a'$)" missing, unbalanced parenthesis

I am also going to introduce a strengthened variant of \textbf{Der} 3 $b$) :

\textbf{Der} 3 $b'$) Let $Z' \to Z$ be a \Hot-equivalence in \Dia. Then $f^* : \DD(Z) \to \DD(Z')$ induces an equivalence between the categories of \underline{locally} \underline{constant} coefficients on $Z$ and on $Z'$.\footnote[50]{\textbf{Der} 3 $b'_0$) : the functor is faithful\\
\hspace*{6em}$b'_1$) : the functor is fully faithful\\
\hspace*{6em}$b'_2$) : the functor is an equivalence} % [sic] localements -> localement, plural ending on the adverb

\msp{page 99}
In certain questions it suffices to know that the said functor is fully faithful. This implies that
\[
H^*_{\DD}(Z,F) \overset{\sim}{\longrightarrow} H^*_{\DD}(Z', f^*F)
\]
when $F$ is \underline{locally} \underline{constant} on $Z$\footnote[51]{One may say that $f : Z' \to Z$ is a $\DD$-equivalence 0-\textit{complete on the left,} and 0-\textit{complete on the right} in the homological case.} (but it is not clear that this condition always implies full faithfulness\footnote[52]{But this is OK if one has internal \underline{\Hom} in $\DD(Z)_{\mathrm{loc.const.}}$ satisfying
\[
\Hom(F,G) \simeq \Hom(e_Z, \underline{\Hom}(F,G)) \quad ,
\]
where $e$ is the final element of $\mathcal{A}$, the \underline{\Hom} being compatible with inverse images.}) and similarly for homology
\[
H_*^{\DD}(Z', f^*F) \longrightarrow H_*^{\DD}(Z,F) \quad .
\]
Under this condition, we have the :

\textbf{Proposition :} \textit{We assume} \textbf{Der} 3 \textup{b}$'$) \textit{holds. Let $f : X \to Y$ in} \Dia. \textit{For $f$ to be a $\DD$-fibration, it is necessary and sufficient that for $Y'' \to Y'$ over $Y$ which is a} \Hot\textit{-equivalence, $X'' = X \times_Y Y'' \to X' = X \times_Y Y'$ be a $\DD$-equivalence.}

To tell the truth, one can state something reasonable, without assuming axiom \textbf{Der} 3 $b'$). One must introduce the :

\textbf{Definition :} Let $f : Z' \to Z$ in \Dia. We say that $f$ and a 1-\textit{complete} $\DD$-equivalence if the functor $f^* : \DD(Z) \to \DD(Z')$ is fully faithful on the full subcategory of $\DD(Z)$ consisting of the locally constant coefficients. We say that $f$ is a 2-\textit{complete} $\DD$-equivalence if it induces an equivalence between the categories of locally constant coefficients. % [sic] f et une -> f est une, verb "est" misprinted as "et"

\msp{page 100}
We say that $Z$ is \textit{completely} $\DD$-\textit{aspheric}, if $Z \to e$ is a 2-complete $\DD$-equivalence.

\printedpage{61}
% head: A. Grothendieck : Dérivateurs, Ch. I, Édition des Universités de Montpellier II et Paris VII
% foot: Version 8 juillet 2001

Thus, axiom \textbf{Der} 3 $a'$) can be reformulated, by saying that if $Z$ is \Hot-aspheric, then $Z$ is completely $\DD$-aspheric. And likewise, \textbf{Der} 3 $b'$) by saying that if $f : X \to Y$ in \Dia{} is a \Hot-equivalence, then $f$ is a 2-complete $\DD$-equivalence. We are going to give a weakened variant, of \textbf{Der} 3 $a'$), which suffices to give a reasonable notion of $\DD$-fibration :

\textbf{Der} 3 $a''$) Every $Z$ in \Dia{} which has either a final object, or an initial object, is completely $\DD$-aspheric.

\textbf{Definition.} We say that $f : X \to Y$ in \Dia{} is a $\DD$-fibration, if for
\[
Y'' \longrightarrow Y' \longrightarrow Y
\]
with $Y'' \to Y'$ a 1-complete $\DD$-equivalence (for example $Y'$ and $Y''$ completely $\DD$-aspheric), $X'' = X \times_Y Y'' \longrightarrow X' = X \times_Y Y'$ is a $\DD$-equivalence. % [sic] asphèriques -> asphériques, grave accent instead of acute

This is equivalent to each of conditions $b$), $b'$) of prop. 2 (page 79), \textit{i.e.} that $f$ be universally a left $\DD$ fibration, or universally a right $\DD$-fibration (thanks to \textbf{Der 3} $b''$ \texttt{[rather} $a''$\texttt{]}). % [sic] $\DD$ fibration -> $\DD$-fibration, hyphen missing

\msp{page 101}
\textbf{Definition 1.} $f : X \to Y$ in \Dia{} is called a \textit{cohomological $\DD$-equivalence} (understood : \textit{complete}) if for every $F \in \Ob\,\DD(Y)$, $F$ \underline{locally} \underline{constant}, we have
\[
H^*_{\DD}(Y,F) \overset{\sim}{\longrightarrow} H^*_{\DD}(X, f^*(F)) \quad .
\]
Dual notion \textit{homological $\DD$-equivalence}
\[
H_*^{\DD}(X, f^*(F)) \overset{\sim}{\longrightarrow} H_*^{\DD}(Y,F) \quad .
\]
If $Y$ is completely $\DD$-aspheric, then the locally constant coefficients on $Y$ are constant, and $f$ is a cohomological (resp. homological) $\DD$-equivalence if and only if it is a $\DD$-equivalence \textit{i.e.} if $X$ is $\DD$-aspheric.\footnote[53]{Axiom \textbf{Der} 3 $b'_1$) implies that every \Hot-equivalence is a cohomological $\DD$-equivalence.}

\textbf{Proposition 1} (Sorital). \textit{Let $X \overset{f}{\to} Y \overset{g}{\to} Z$. If $f$ and $g$ are cohomological $\DD$-equivalences, then $gf$ too. If $gf$ and $f$ are cohomological $\DD$-equivalences, then $g$ too. If $gf$ is a cohomological $\DD$-equivalence, and if $g$ is a \underline{2-complete} $\DD$-equivalence} (hence every $F$ locally constant on $Y$ comes from a $G$ locally constant on $Z$), \textit{then $f$ is a cohomological $\DD$-equivalence.}

\textbf{Definition 2.} 1°) $f : X \to Y$ in \Dia{} is called a \textit{complete cohomological $\DD$-fibration,} I do mean \underline{complete}, if

$a$) For every $y \in \Ob\, Y$, $X_y \to X/y$ is a cohomological $\DD$-equivalence (complete).

$b$) For every $y' \to y$ in $Y$, $X/y' \to X/y$ is a cohomological $\DD$-equivalence (complete).

\printedpage{62}
% head: A. Grothendieck : Dérivateurs, Ch. I, Édition des Universités de Montpellier II et Paris VII
% foot: Version 8 juillet 2001

\msp{page 102}
$c$) These conditions remain satisfied after every base change $Y' \to Y$.

2°) Dual notion of complete homological fibration.

\textsc{Example.} If $f$ is \Hot-proper, the arrows of $a$) and $b$) are \Hot-equivalences, hence by virtue of \textbf{Der 3} $b'_1$) (which must be admitted), they are cohomological $\DD$-equivalences. This means that $a$) and $b$) hold, and $c$) too, for the notion of a \Hot-proper arrow is stable under base change.

\textbf{Proposition 2.} \textit{Let $f : X \to Y$ be a complete cohomological $\DD$-fibration, $F \in \Ob\,\DD(X)$ locally constant on $X$. Then $f_*(F)$ is locally constant, its formation commutes with every base change, and in particular is computed fiber by fiber.}\footnote[54]{There is a converse \textit{cf.} prop. 3$'$.}

\textbf{Corollary 1.} \textit{Let there be a cartesian diagram}
\[
\begin{tikzcd}
X & X' \arrow[l, "q"'] \arrow[d, "f'"] \\
Y & Y' \arrow[l, "p"']
\arrow[from=1-1, to=2-1, "f"']
\end{tikzcd}
\quad .
\]
\textit{We assume $f$ a cohomological $\DD$-fibration and $p$ a cohomological $\DD$-equivalence. Then $q$ is a cohomological $\DD$-equivalence.}

\msp{page 103}
Indeed, we have for $F$ locally constant on $X$
\[
H^*_{\DD}(X,F) \simeq H^*_{\DD}(Y, f_*F) \overset{\sim}{\longrightarrow} H^*_{\DD}(Y', p^*f_*F) \simeq H^*_{\DD}(Y', f'_*q^*F) \simeq H^*_{\DD}(X', q^*F)
\]
qed.

\textbf{Proposition 3.} \textit{The notion of complete cohomological $\DD$-fibration is stable under base change} (trivial) \textit{and under composition.}

For stability under composition, it suffices to note the following corollary of proposition 2.

\textbf{Corollary 2} (of prop. 2). \textit{Let $f : X \to Y$. For $f$ to be a complete cohomological $\DD$-fibration, it is necessary and sufficient that for every $Y'' \overset{p}{\to} Y' \to Y$ over $Y$, with $p$ a cohomological $\DD$-equivalence, the corresponding arrow $X'' = X \times_Y Y'' \longrightarrow X' = X \times_Y Y'$ also be a cohomological $\DD$-equivalence.}

This is necessary by corollary 1, taking into account that $f' : X' = X \times_Y Y' \to Y'$ is also a cohomological $\DD$-fibration. It is sufficient, for it suffices to apply it to the cases $\{y\} \to Y/\{y\}$, $Y/y \to Y/y'$, which are cohomological $\DD$-equivalences by
\msp{page 104}
virtue of \textbf{Der} 3 $b'_1$) (or only \textbf{Der} 3 $a'$)).

\printedpage{63}
% head: A. Grothendieck : Dérivateurs, Ch. I, Édition des Universités de Montpellier II et Paris VII
% foot: Version 8 juillet 2001

\textbf{Proposition 4.} \textit{We assume that} \textbf{Der} 3 a$'$) \textit{is satisfied} (\textit{i.e.} every $Z$ in \Dia{} which is \Hot-aspheric, is completely $\DD$-aspheric). \textit{Let $f : X \to Y$ in} \Dia. \textit{The following conditions are equivalent}

1°) \textit{$f$ is a \underline{complete} cohomological $\DD$-fibration} (conditions $a$), $b$), $c$) of \texttt{[the]} definition 2, page 101).

2°) \textit{For every $F \in$} \Ob\,$\DD(X)$\textit{, locally constant, $f_*(F)$ is locally constant, is computed fiber by fiber; and that remains true after every base change $Y' \to Y$ (with $F'$ locally constant on $X' = X \times_Y Y'$). (Then the formation of $f_*(F)$ commutes with \underline{every} base change).}

3°) \textit{For every $Y'' \overset{p}{\to} Y'$ over $Y$, if $p$ is a cohomological $\DD$-equivalence, then $X'' = X \times_Y Y'' \longrightarrow X' = X \times_Y Y'$ too.}

3$'$) \textit{As in 3, assuming that $Y''$ is the one-point category, and that $Y'$ has a final element : $X'' \to X'$ must be a cohomological $\DD$-equivalence.}

\msp{page 105}
\textbf{Definition 3.} We say that $f : X \to Y$ in \Dia{} is a \textit{complete $\DD$-fibration} if it is a complete cohomological $\DD$-fibration and a complete homological $\DD$-fibration.

This is stable under base change, under composition.

Definition 3 is equivalent to this : for every $p : Y'' \to Y'$ over $Y$, with $Y''$ and $Y'$ \Hot-aspheric, the corresponding morphism $X'' \to X'$ is \underline{at} \underline{the same} \underline{time} a cohomological $\DD$-equivalence, and a homological $\DD$-equivalence. (And it even suffices to assume it for $Y'' = e$ and for $Y'$ with a final object or an initial object.) One can think of strengthening this condition still further by requiring (without looking to the left or to the right) that $X'' \to X'$ be a 1-complete $\DD$-equivalence. But it is not clear that this strengthened notion is stable under composition – unless one strengthens it still further, by requiring that if $Y'' \to Y'$ is a 1-complete $\DD$-equivalence, then so is $X'' \to X'$. The trouble with this strengthened definition is that for lack of a cohomological or homological description of the notion of
\msp{page 106}
1-complete $\DD$-equivalence, I do not feel capable of giving interesting (not totally trivial) examples of $\DD$-fibrations in this very strong sense. Except for the trivial \Hot-fibrations (\textit{i.e.} with \Hot-aspheric fibers) – but that is no fun!

But when one has associative and unital products $*$ and corresponding internal \underline{\Hom} in the $\DD(X)_{\mathrm{loc.\ const.}}$, compatible with \texttt{[the]} inverse images, so one will have on $X$ if $F$ and $G$ are locally constant in $\DD(X)$, and if $\varepsilon_X$ denotes the unit for $*$ in $\DD(X)$,
\[
\Hom(F,G) \simeq \Hom(\varepsilon_X * F, G) \simeq \Hom(\varepsilon_X, \underline{\Hom}(F,G)) \simeq \Hom_{\mathcal{A}}(\varepsilon_e, p_*\, \underline{\Hom}(F,G)) \ ,
\]
for $\varepsilon_X = p^*(\varepsilon_e)$, where $p : X \to e$, so in this case an arrow $f : X \to Y$ in \Dia{} is a 1-complete $\DD$-equivalence if (and only if) it is a cohomological $\DD$-equivalence. So in this case, if $f$ is a complete cohomological $\DD$-fibration, it is also a complete homological $\DD$-fibration, hence also a complete $\DD$-fibration (definition 3). (And even in the so-called strengthened sense of just now : if $Y'' \to Y'$ over $Y$ is a 1-complete $\DD$-equivalence ($\Leftrightarrow$ a cohomological $\DD$-equivalence!) then $X'' \to X'$ too. % [sic] soit-disant -> soi-disant, spelling; % [sic] "(Et même ... aussi." -> closing parenthesis missing, unbalanced parenthesis

\end{document}
